\section{时变线性化模型预测控制}
\subsection{参考与离散误差模型}
从通过验收的纯PSO结果中选择复核目标最小者作为参考，不施加任何局部精修。参考包含$\bar x(t)=(r,\theta,u,w,m)^{\mathsf T}$和$\bar\delta(t),\bar q(t)$；预测窗口超出计划关机时刻后，参考继续按无推力、无随机外力的动力学延伸。

采样周期$T_s=2$\,s、预测步数$N=15$。本周期施加$\delta(t)=\bar\delta(t)+\Delta\delta_k$、$q(t)=\bar q(t)+\Delta q_k$，其中修正$c_k=(\Delta\delta_k,\Delta q_k)^{\mathsf T}$保持常数，参考前馈仍连续变化。以RK4积分名义方程得到离散映射$F_k(x,c)$，沿参考展开并归一化：
\begin{equation}e_{k+1}=A_ke_k+B_kv_k+d_k,\label{eq:ltv}\end{equation}
\begin{equation}
\begin{aligned}
e_k&=S_x^{-1}(x_k-\bar x_k),&v_k&=S_c^{-1}c_k,\\
A_k&=S_x^{-1}F_{x,k}S_x,&B_k&=S_x^{-1}F_{c,k}S_c,\\
d_k&=S_x^{-1}[F_k(\bar x_k,0)-\bar x_{k+1}].
\end{aligned}
\end{equation}
$F_{x,k},F_{c,k}$由离散映射中心差分获得。状态尺度为$(500\,\mathrm m,500\,\mathrm m/R_E,30\,\mathrm{m/s},30\,\mathrm{m/s},1000\,\mathrm{kg})$，输入尺度为$(2^\circ,50\,\mathrm{kg/s})$。$d_k$保留参考插值、积分和名义扰动环境差异，不能解释为未来实际噪声预测。名义预测不读取实际未来扰动。

\subsection{代价凝聚与凸二次规划}
未来状态及控制堆叠为$E,U$，则式\eqref{eq:ltv}递推得到
\begin{equation}E=\mathcal A e_k+\mathcal B U+\mathcal D=a+\mathcal B U.\end{equation}
取有限时域目标
\begin{equation}
\begin{aligned}
J_{\rm MPC}={}&E^{\mathsf T}\bar Q E+U^{\mathsf T}\bar R U\\
&+(DU-b_\Delta)^{\mathsf T}\bar R_\Delta(DU-b_\Delta)
+\lambda_\epsilon\epsilon^{\mathsf T}\epsilon .
\end{aligned}
\end{equation}
$D$为控制差分矩阵，$b_\Delta$第一块为上步已执行修正。状态权重$Q=\operatorname{diag}(4,2,4,4,0.1)$，末步权重为$5Q$，$R=0.2I_2$，$R_\Delta=0.5I_2$，$\lambda_\epsilon=10^5$。惩罚的是反馈修正变化，不是总控制的硬性变化率限制。

令$\zeta=(U^{\mathsf T},\epsilon^{\mathsf T})^{\mathsf T}$，消去状态并去除常数项，得到
\begin{equation}\min_\zeta\tfrac12\zeta^{\mathsf T}H\zeta+f^{\mathsf T}\zeta,\quad G\zeta\le b,\quad l\le\zeta\le u,\end{equation}
其中
\begin{equation}
\begin{aligned}
H&=\operatorname{blkdiag}\left(2[\mathcal B^{\mathsf T}\bar Q\mathcal B+\bar R+D^{\mathsf T}\bar R_\Delta D],2\lambda_\epsilon I_N\right),\\
f&=\begin{bmatrix}2(\mathcal B^{\mathsf T}\bar Q a-D^{\mathsf T}\bar R_\Delta b_\Delta)\\0\end{bmatrix}.
\end{aligned}
\end{equation}
因$\bar R\succ0$、$\bar Q,\bar R_\Delta\succeq0$，目标严格凸\cite{boyd2004}。状态选择矩阵$C_r,C_m$分别提取径向及质量误差；参考高度和燃料余量归一化为$g,b_m$后，
\begin{equation}
G=\begin{bmatrix}-C_r\mathcal B&-I_N\\-C_m\mathcal B&0\end{bmatrix},\quad
b=\begin{bmatrix}g+C_ra\\b_m+C_ma\end{bmatrix},\quad\epsilon\ge0.
\end{equation}
地面为软约束，干质量为硬约束；输入限制$|\delta|\le15^\circ$、$0\le q\le350$\,kg/s，修正幅度限制$|\Delta\delta|\le5^\circ$、$|\Delta q|\le100$\,kg/s。滑行及关机后不新增点火。当前QP含$2N+N=45$个变量，以凸内点法求解。

每次仅执行第一项修正，下一周期更新状态后重算。起飞后首次反馈在6\,s施加。假设理想全状态测量，不包含传感器或执行机构动态；QP失败时退回参考前馈。该回退不构成一般安全性保证，局部QP凸性也不意味着原非线性控制问题具有全局最优性或递归可行性。
